%====================================================%
% SEZIONE 1: Teoria
%====================================================%
\begin{question}[ID=1.1, subtitle={Domande di teoria sulle disequazioni}, difficulty=facile, tags={teoria, disequazioni}]
Rispondi alle seguenti domande di teoria:
\begin{center}\begin{minipage}{0.85\linewidth}
\begin{enumerate}[label=\arabic*)]
\item Che cos'è una disequazione di primo grado?
\item Qual è il primo principio di equivalenza delle disequazioni?
\item Qual è il secondo principio di equivalenza delle disequazioni?
\item Come si rappresenta graficamente una soluzione su una retta?
\item Quando una disequazione è sempre verificata?
\item Quando una disequazione è impossibile?
\item Quali operazioni si possono fare mantenendo l'equivalenza?
\item Cosa succede al verso moltiplicando per un numero positivo?
\item Cosa succede al verso moltiplicando per un numero negativo?
\item Quando bisogna cambiare il verso della disequazione?
\item Cos'è una disequazione frazionaria?
\item Qual è la condizione di esistenza in una disequazione frazionaria?
\item Perché è importante studiare il segno del denominatore?
\item Come si risolve una disequazione con frazioni?
\item Cos'è il metodo del segno?
\item Come si costruisce una tabella dei segni?
\item Cosa sono le C.E.?
\item Perché il denominatore non può annullarsi?
\item Come si scrive la soluzione in forma di intervallo?
\item Qual è la differenza tra intervallo aperto e chiuso?
\end{enumerate}
\end{minipage}\end{center}
\begin{solution}Risposta orale o scritta per ciascun punto.\end{solution}
\end{question}

%====================================================%
% SEZIONE 2: Disequazioni lineari
%====================================================%
\begin{question}[ID=2.1, subtitle={Disequazioni di primo grado}, difficulty=facile, tags={disequazioni lineari}]
Risolvi le seguenti disequazioni di primo grado:
\begin{center}\begin{minipage}{0.92\linewidth}
\begin{multicols}{2}
\begin{enumerate}[label=\arabic*),leftmargin=3.2cm]
\item $x + 3 > 7$
\item $x - 5 \leq 2$
\item $2x > 8$
\item $-3x \leq 9$
\item $2x + 1 > 5$
\item $3x - 4 \geq 2$
\item $5 - x < 3$
\item $4 - 2x \geq 6$
\item $2(x + 1) > 6$
\item $3(x - 2) \leq 6$
\item $2x + 3 > x - 1$
\item $4x - 5 \leq 2x + 3$
\item $3x - 2 > 5x - 6$
\item $2 - x \geq x + 4$
\item $\frac{x}{2} > 3$
\item $\frac{3x - 1}{2} \leq 4$
\item $\frac{x + 2}{3} > \frac{x - 1}{2}$
\item $\frac{2x - 3}{4} \leq \frac{x + 5}{2}$
\item $3(2x - 1) - 4x > 5$
\item $2(3x + 1) - (x - 4) \leq 7$
\end{enumerate}
\end{multicols}
\end{minipage}\end{center}
\begin{solution}
1) $x>4$; 2) $x\le 7$; 3) $x>4$; 4) $x\ge -3$; 5) $x>2$; 6) $x\ge 2$; 7) $x>2$; 8) $x\le -1$; 9) $x>2$; 10) $x\le 4$; 11) $x>-4$; 12) $x\le 4$; 13) $x<2$; 14) $x\le -1$; 15) $x>6$; 16) $x\le 3$; 17) $x<7$; 18) $x\in\mathbb{R}$; 19) $x>4$; 20) $x\le\tfrac{1}{5}$
\end{solution}
\end{question}

%====================================================%
% SEZIONE 3: Disequazioni frazionarie
%====================================================%
\begin{question}[ID=3.1, subtitle={Disequazioni frazionarie}, difficulty=media, tags={disequazioni frazionarie}]
Risolvi le seguenti disequazioni frazionarie:
\begin{center}\begin{minipage}{0.92\linewidth}
\begin{multicols}{2}
\begin{enumerate}[label=\arabic*),leftmargin=3.2cm]
\item $\frac{x}{x-1} > 0$
\item $\frac{x-2}{x+1} \ge 0$
\item $\frac{2x+1}{x-3} > 0$
\item $\frac{x+4}{2x-1} \le 0$
\item $\frac{x-1}{x+2} < 0$
\item $\frac{3x-2}{x-2} \ge 0$
\item $\frac{x+1}{x-4} > 1$
\item $\frac{2x}{x-3} \le 1$
\item $\frac{x-2}{x+1} > \frac{x}{x+1}$
\item $\frac{2x+3}{x-1} \ge \frac{x}{x-1}$
\item $\frac{x+1}{x-2} < \frac{2x}{x-2}$
\item $\frac{3x-1}{x+2} \le \frac{x+4}{x+2}$
\item $\frac{x}{3x-1} > 1$
\item $\frac{2}{x+1} \ge 2$
\item $\frac{x+1}{x-2} - 1 > 0$
\item $\frac{2x}{x+3} - 2 \le 0$
\item $\frac{x^2-1}{x-1} > 0$
\item $\frac{x^2-4}{x+2} \le 0$
\item $\frac{3}{x-2} + \frac{1}{2x+1} > 0$
\item $\frac{1}{x+1} - \frac{3}{x-1} \le 0$
\end{enumerate}
\end{multicols}
\end{minipage}\end{center}
\begin{solution}
1) $x<0\vee x>1$; 2) $x\le -1$ escluso $\cup x\ge 2$; 3) $x<-\tfrac12\vee x>3$; 4) $-4\le x<\tfrac12$; 5) $-2<x<1$; 6) $x\le\tfrac23\vee x>2$; 7) $4<x<5$... vedi svolgimento completo.
\end{solution}
\end{question}
